\begin{afterword}
\summaryhead

Blood types are real categories. Landsteiner's discovery of them made transfusion safe, and a mismatch between mother and child can endanger a second child. Yet in Japan blood types carry personality stereotypes, as astrological signs do in the West, down to horoscopes in the newspaper. Because blood types were never mysterious and never the lines of a conflict, the author argues, the stereotypes must have arisen from the mere existence of the labels.

The author does not recall an experiment that labels plants, rocks or furniture, but predicts that labeling would have the same power over them. Four explanations are offered: a boundary makes the mind look for similarities inside it; a name sets aside a neural subnetwork; things with one name are compressed into one point; and people make up claims about any named category. A Western blood-type diet book was a bestseller. The post concludes that drawing a category is not a neutral act for a human mind, which is one more reason not to believe you can define a word any way you like.

\responsehead

The post cites no studies. Its checkable claims are facts about blood, a history, and a prediction about experiments.

The facts about blood are correct. Landsteiner won the 1930 Nobel Prize for discovering human blood groups, and Landsteiner also identified the Rh factor. Japanese newspapers do print blood type horoscopes, and the diet book was a \textsc{New York Times} bestseller.

The main inference is stronger than the history allows. The post says the stereotypes ``must have arisen strictly from the mere existence of the labels.'' As Wikipedia summarizes the history, they have authors: Takeji Furukawa, whose 1927 paper linked blood type and temperament and who later explained Taiwanese resistance to Japanese rule by the share of type O blood, and Masahiko Nomi, whose books of the 1970s are where the popular belief comes from. The case shows people believing claims that someone wrote about a category, which is the fourth of the post's explanations. It does not show labels alone producing stereotypes, which is what the post uses it for.

The experiment the author could not recall had been done, more than once, before the post. Carmichael, Hogan and Walter (1932) found that naming a drawing just before it was shown could change how people later reproduced it. Tajfel and Wilkes (1963) found that lines labeled by length were judged farther apart. These support the prediction. The second study also found a limit: a classification ``not coherently related'' to length did not have this effect. On that evidence, labels change judgement mainly when they go with what is judged, which qualifies ``mere labeling had power over all things.''

Finally, the conclusion is stated for all categories: ``as soon as you actually think of them, they exert force on your mind.'' The post's own evidence is one category of people and a prediction for others.

In short: accurate facts about blood types, a history of the stereotypes that points to their authors more than to the labels, and a prediction that earlier experiments had already tested and supported, with a limit on ``mere'' labeling.
\end{afterword}
